% Unite: papier 34, section 11; traduction francaise depuis la source allemande auditee.

\providecommand{\mo}{\mathfrak{o}}
\providecommand{\ma}{\mathfrak{a}}
\providecommand{\mA}{\mathfrak{A}}
\providecommand{\mZ}{\mathfrak{Z}}

\subsection*{\S{} 11. Le centre}

Le centre $\mZ$ d'un anneau $\mo$ est l'ensemble des $z$ qui commutent avec tous les éléments de l'anneau ($za=az$ pour tout $a$). $\mZ$ est un anneau commutatif.

Tout idéal unilatère $\ma$ dans $\mo$ possède dans $\mZ$ un ``idéal contracté'' $\ma\cap\mZ$. En effet, puisque $\ma$ aussi bien que $\mZ$ sont des $\mZ$-modules, $\ma\cap\mZ$ l'est aussi.

Tout idéal $\mA$ dans $\mZ$ possède un idéal étendu: l'idéal $\ma=(\mA,\mo\mA)$ engendré par $\mA$ dans $\mo$. Celui-ci est bilatère, car
\[
        \ma\mo=(\mo\mA,\mA)\mo=(\mo\mA\mo,\mA\mo)
        = (\mo^2\mA,\mo\mA)\subseteq \mo\mA\subseteq\ma.
\]
Si l'on suppose donné dans $\mo$ un élément unité, on peut écrire: $\ma=\mo\mA$. Dans ce cas vaut le théorème suivant:

\emph{À toute décomposition bilatère de $\mo$ correspond biunivoquement une décomposition du centre: de}
\begin{equation}
        \mo=\ma_1+\cdots+\ma_n,
        \qquad \mA_i=\ma_i\cap\mZ                         \tag{1}
\end{equation}
\emph{on obtient}
\begin{equation}
        \mZ=\mA_1+\cdots+\mA_n,
        \qquad \mo\mA_i=\ma_i,                              \tag{2}
\end{equation}
\emph{et réciproquement.}

\emph{Démonstration.} Soit $z$ un élément du centre, et
\begin{equation}
        z=z_1+\cdots+z_n                                    \tag{3}
\end{equation}
sa décomposition suivant (1). Alors les $z_i$ sont des éléments du centre; car
\[
        za=z_1a+\cdots+z_na=az=az_1+\cdots+az_n;
\]
à cause du caractère direct de la somme, et puisque les $\ma_i$ sont bilatères,
\[
        z_i a=az_i.
\]
Donc $z_i$ appartient à $\ma_i\cap\mZ=\mA_i$; ainsi (3) donne la décomposition en somme affirmée. En particulier on obtient $e=\sum e_i$; donc les $e_i$ sont des éléments du centre. Enfin
\[
        z=ze_1+\cdots+ze_n,
\]
donc
\[
        \mA_i=\mZ e_i,
        \qquad
        \mo\mA_i=\mo\mZ e_i=\mo e_i=\ma_i.
\]

Réciproquement, si $\mZ=\mA_1+\cdots+\mA_n$ est une décomposition du centre, alors, d'après le \S{}9,
\[
        \mA_i=\mZ e_i,
        \qquad e_i^2=e_i,
        \qquad e_i e_k=0\quad(i\ne k).
\]
Il vient donc, à cause des relations d'orthogonalité \hbox{\(\S 9\)},
\[
\begin{aligned}
        \mo&=\mo e=\mo e_1+\cdots+\mo e_n
        =\mo\mZ e_1+\cdots+\mo\mZ e_n\\
        &=\mo\mA_1+\cdots+\mo\mA_n
        =\ma_1+\cdots+\ma_n .
\end{aligned}
\]
Si l'on pose maintenant $\ma_i\cap\mZ=\mA_i'$, alors on sait, par la première partie de l'assertion:
\[
        \mZ=\mA_1'+\cdots+\mA_n'\quad\hbox{directement}.
\]
Mais $\mA_i'\supseteq\mA_i$, donc $\mA_i'=\mA_i$ (\S 4, 2).

\emph{Corollaire.} $\ma_i$ et $\mA_i$ sont simultanément décomposables directement en idéaux bilatères, ou ne le sont pas.
